% 题干转写：完成
% 解析转写：完成
% 图形核对：完成
\section{第35届复赛第2题：弹簧--摩擦阻尼往返运动}
\subsection{题干}
如图\ref{fig:cpho35q02}，劲度系数为 $k$ 的轻弹簧左端固定、右端连一质量为 $m$ 的小球。弹簧自然状态时小球位于原点 $O$；小球在粗糙水平面上滑动，动摩擦因数为 $\mu$，最大静摩擦力等于滑动摩擦力。初始时小球静止，弹簧伸长为 $-A_0$，其中 $A_0>0$，重力加速度为 $g$。
\begin{figure}[htbp]\centering
\begin{tikzpicture}[>=Stealth,scale=0.9]
  \draw[thick] (-4.4,0) -- (4.6,0);
  \foreach \x in {-4.2,-3.9,...,4.4} {\draw[gray] (\x,0) -- +(0.18,-0.13);}
  \draw[->,thick] (-4.1,0.36) -- (4.25,0.36) node[right] {$x$};
  \draw[thick] (-3.8,-0.55) -- (-3.8,1.25);
  \foreach \y in {-0.45,-0.15,...,1.15} {\draw[thick] (-4.08,\y) -- (-3.8,\y+0.15);}
  \draw[thick,decorate,decoration={coil,segment length=4pt,amplitude=3pt}]
    (-3.8,0.65) -- (-1.15,0.65);
  \fill[blue!25] (-0.82,0.33) rectangle (-0.17,0.98);
  \draw[thick] (-0.82,0.33) rectangle (-0.17,0.98);
  \node at (-0.50,0.66) {$m$};
  \draw[densely dashed] (-0.50,0.1) -- (-0.50,1.1);
  \draw[densely dashed] (1.15,0.1) -- (1.15,1.1);
  \node[below] at (-0.50,0.1) {$-A_0$};
  \node[below] at (1.15,0.1) {$O$};
  \node[above] at (-0.50,1.08) {$t=0$};
  \node[below] at (-2.45,-0.5) {劲度系数 $k$ 的轻弹簧};
\end{tikzpicture}

\caption{第35届复赛第2题题图的TikZ重绘：初始压缩状态、原点和水平坐标轴。}\label{fig:cpho35q02}
\end{figure}
\begin{enumerate}
\item 小球至多只能向右运动时，求最终位置及 $A_0$ 条件；
\item 小球第一次向右越过原点后至多只能向左运动时，求最终位置及条件；
\item 小球只能完成 $n$ 次往返运动时，求最终位置及条件；
\item 在只能完成 $n$ 次往返运动时，求全过程总路程。
\end{enumerate}

\subsection{参考解答整理}
记 $q=\mu mg/k$。每一次由静止位置到下一静止位置的过程中，功能原理给出弹簧势能减少等于摩擦耗散。

\paragraph{（1）至多向右运动。} 若 $0<A_0\leq q$，初始弹力不足以克服静摩擦，小球停在 $x=-A_0$。若 $A_0>q$，第一次向右停下的位置为
\[x=A_0-2q.\]
因此在“至多向右运动”的范围内，最终位置分段为
\[
x=\begin{cases}
-A_0,&0<A_0\leq q,\\
A_0-2q,&q<A_0\leq2q,\\
A_0-2q,&2q<A_0\leq3q.
\end{cases}
\]
第二、三段分别落在原点左侧（或原点）和右侧；其后均因回复弹力不超过 $\mu mg$ 而静止。

\paragraph{（2）至多向左运动。} 当 $3q<A_0\leq4q$ 时，第一次向左停下后仍在原点右侧或原点：
\[x=4q-A_0\geq0.\]
当 $4q<A_0\leq5q$ 时，最终停在原点左侧：
\[x=4q-A_0=-A_1<0.\]

\paragraph{（3）完成 $n$ 次往返后的最终位置。} 令第 $n$ 次返回原点左侧的静止位置为 $-A_n$。相邻转折点的伸长量依次相差 $2q$，故
\[A_n=A_0-4nq.\]
最终位置的完整分段为
\[
x=\begin{cases}
4nq-A_0,&4nq<A_0\leq(4n+1)q,\\
A_0-2(2n+1)q,&(4n+1)q<A_0\leq2(2n+1)q,\\
A_0-2(2n+1)q,&2(2n+1)q<A_0\leq(4n+3)q,\\
4(n+1)q-A_0,&(4n+3)q<A_0\leq4(n+1)q.
\end{cases}
\]
其中第二、四段的右端点分别包含原点情形。

\paragraph{（4）总路程。} 从初态到最终静止位置 $x$ 的能量关系为
\[\frac12kA_0^2-\frac12kx^2=\mu mgs,\qquad s=\frac{k}{2\mu mg}(A_0^2-x^2).\]
代入上述分段终点，得到
\[
s=\begin{cases}
4n(A_0-2nq),&4nq<A_0\leq(4n+1)q,\\
2(2n+1)\bigl[A_0-(2n+1)q\bigr],&(4n+1)q<A_0\leq(4n+3)q,\\
4(n+1)\bigl[A_0-2(n+1)q\bigr],&(4n+3)q<A_0\leq4(n+1)q.
\end{cases}
\]

\subsection{转写依据}
题干、答案分段和题图分别依据 \texttt{CPHO35-Q02.md}、\texttt{CPHO35-Q02\_formula\_check\_v1.csv} 及图件核对入口。分值边界未在当前PDF中明示，未写入本条目。
